%贵州大学选修课课程报告LaTeX版（非官方）
%请使用XeLaTeX进行编译（若需要目录，请编译两次）
%版本：1.0
%作者：Spico(spico@qq.com)

\documentclass[12pt,a4paper]{ctexart} %使用CTeX文档包，正文部分字号12，A4纸张
\usepackage{ulem}   %用于字体加强
\usepackage{amsmath}    %用于数学公式
\usepackage{graphicx}   %用于插入图像
\usepackage{subfig} %用于排列图像
\usepackage{booktabs}   %用于制作三线表格
\usepackage{multirow}   %用于表格跨行
\usepackage{fancyhdr}   %用于设置页面格式
\usepackage{verbatim}   %用于打印代码（原文输出）
\usepackage{setspace}   %设置行距

\pagestyle{fancy}
\fancyhf{}
\lhead{\textbf{课程报告}}
\rhead{\thepage}
\cfoot{}
%\setCJKmainfont{宋体}
\newcommand{\Date}[3]{\heiti\zihao{2} #1年#2月#3日}    %生成日期

\begin{document}

\centering
%第一页：封面
\begin{figure}[h]
    \centering
    \subfloat{
        \label{fig:schoolBadge}
        \includegraphics[keepaspectratio,width=90pt]{pic/school_badge.png}
    }
    \hspace{10pt}
    \subfloat{
        \label{fig:schoolName}
        \includegraphics[keepaspectratio,width=200pt]{pic/school_name.png}
    }
    \newline\newline
    \centering
    \heiti\zihao{-0}　贵州大学课程报告 \newline\newline\newline
    \heiti\zihao{1}　中国酒概述---酒文化之我见 \newline\newline
\end{figure}

%生成署名信息
\begin{tabular}{rcl}
    \centering
    \zihao{2}\heiti{学\quad 院：} & \zihao{2}\uline{\makebox[9em][c]{\fangsong 大数据与信息工程}} \\ \\
    \zihao{2}\heiti{专\quad 业：} & \zihao{2}\uline{\makebox[9em][c]{\fangsong 通信工程}} \\ \\
    \zihao{2}\heiti{班\quad 级：} & \zihao{2}\uline{\makebox[9em][c]{\fangsong 通信152班}} \\ \\
    \zihao{2}\heiti{学\quad 号：} & \zihao{2}\uline{\makebox[9em][c]{\fangsong 15********}} \\ \\
    \zihao{2}\heiti{姓\quad 名：} & \zihao{2}\uline{\makebox[9em][c]{\fangsong 太子夜华}} \\ \\ \newline\newline\newline
\end{tabular}

\Date{2017}{2}{19}
%设置页面格式
\setcounter{page}{0}    %设置页码为0
\thispagestyle{empty}   %清除页眉页脚
\newpage    %强制换页

%目录页
\zihao{3}\textnormal{
\setcounter{tocdepth}{2}    %设置目录深度，生成目录
\tableofcontents}    %有目录需要编译两次
\setcounter{page}{0}
\thispagestyle{empty}
\newpage

%正文部分
\flushleft
\section{目的及意义}
\subsection{目的}
\paragraph{\textnormal{　　提起中国白酒，不得不联想起我们的国酒——茅台。茅台，作为酱香型白酒的典范之作，以其酱香突出、香气扑鼻的特点深受国内外饮酒人士的好评。以茅台为代表的贵州名优白酒作为一张贵州旅游的名片，而逐渐为大家所熟知。那么如何充分利用这一酒文化的名片，来塑造贵州良好的旅游形象？本篇将从酒文化之于旅游事业发展的目的、意义、酒文化的发展、促进关系等多方面来论述这一观点。}}
\subsection{意义}
\paragraph{\textnormal{　　自2013年起，我省经济增速就跃居全国前列，甚至达到了增速第一省份的位置。面对这样的经济转型时期，我们要把握好机遇，加速产业链的形成，以达到更好的增长效果，谋求可持续发展，规避泡沫经济。形成良好的经济增长态势，以此提高人民幸福感。
\newline 　　酒文化和贵州息息相关。以茅台酒为例，其历史可以追溯至公元前135年的西汉汉武帝时期\cite{history}。溯本逐源，我们有着悠久的酒文化历史，随着几千年的传承，这一制酒的古法工艺更是被发挥到极致。如果我们可以充分发挥自身优势，将旅游和我们传统的酿酒、饮酒文化相结合，必定能促进整个贵州旅游产业链的发展，实现利益最大化。提升贵州形象，也带领大家走进贵州，充分了解贵州，真正实现“醉美多彩贵州”的文化弘扬。}}
\section{酒文化的内容及其发展}
\paragraph{\textnormal{酒文化是包括酒艺、酒德、酒俗、酒文艺、酒建筑在内的与酒有关的文化体系\cite{culture}。\newline 　　 酒文化即是物质需要与精神需要达到满足的统一。人们丰收而酿造，祭祀以求安。酒自古以来就被赋予了多重含义，成为情感的载体。例如，“浊酒一杯家万里，燕然未勒归无计”，通过杯酒之言生动地描绘出戍边将士心中的愤慨和无奈。又如，“明月几时有，把酒问青天”，每逢佳节倍思亲，古代文人普遍将这种思念寓于酒上。酒承载着人们太多强烈的情感。每感心力焦急，古人便饮酒作赋，以抒发心中的怅然之志。“酒不醉人人自醉”，酒只是一种含有酒精的饮品，它却被历代文人墨客所推崇，以致衍生出如此众多的情感，无不彰显着它源远流长的历史厚重感。人生不如意之事十之八九，每每迁谪异地，总是要取出美酒，痛饮几杯，沉沉睡去，管他这纷繁尘世，管他那昏庸无能，心中的感情，全融在这一杯酒之中，咽了，化了，随他去了。身为千里马，干的却是拉磨的活计，没有伯乐来识别我，罢了罢了，“上善若水，水善利万物而不争”，唯有这汤汤水水懂我的心，和着酒精的麻痹，让我暂时忘记这世间，只蜷缩在美好的梦里，我便无欲无求。
除了感情含义以外，酒文化还融于制酒酿酒的技艺之中。我省在1993-1997年间于仁怀修建了中国酒文化城，意在发扬这种传统的酒文化特色，传播酒文化知识。可以说，酒文化即是一种时代文化。这其中包含了政治、经济、民俗等多方面的文化。例如，我们可以从酒器之中 \cite{tool} 窥探一个朝代的文化。唐代的酒器华丽而精美，反映出这一时代特点，即是国力昌盛，饮酒之风盛行。杜甫在人生的最低谷还不忘吟一句“潦倒新停浊酒杯”，其对酒的痴迷和偏爱程度可见一斑。也反映出某些文人在酒中培养飘然若仙的审美情趣，以成佳句的时代特点。}}
\section{贵州酒文化对本土旅游业的促进作用}
\paragraph{\textnormal{　　既然贵州有如此源远流长的酒历史、酒文化，那么我们怎么充分利用这一优势，使其促进我省旅游业的发展呢？\newline 　　贵州，作为西南的重要省份，其独特的少数民族文化一直吸引着众多游客前来观赏。可以说，贵州拥有众多的旅游资源。但是，我们是否真正地把名片打了出去？很多人提起贵州可能第一印象是黄果树，之后是茅台。那么既然茅台能作为一种贵州的旅游名片而被人们所广为熟知，我们何不扩大这一文化的影响力，借此良机提升自身形象，真正地把这张名片打出去呢？\newline 　　 从功能方面来说，酒文化之于旅游业主要有四方面的功能，即：体验功能、娱乐功能、审美功能\cite{function}。\newline 　　 首先，旅游，即是游客来到实地参观游玩，那么他们就自然而然地向景点提出更高的要求：要有体验的效果，让游客亲身参与其中，才能是好的旅游项目。例如，我们参观某个历史古迹时，总会情不自禁的回想起当年的一些名人轶事。站在故宫的金銮殿之内，不自然地就会回想起明清皇帝朝政的情形；站在岳阳楼，不自然地会想起范仲淹的那句“先天下之忧而忧，后天天下之乐而乐”这些都归功于场景的代入感。可以说，这种代入感在我们的酒文化背景下是比较容易建立的。它需要比较深厚的历史背景和一系列的民俗活动，让游客亲身参与其中。比如前文中提到的水族朋友们的拦门酒，这就给笔者强烈的文化代入感。\newline 　　其次，它还具备娱乐和审美功能。众所周知，我们外出游玩是为了让自己疲惫的身心真正放松下来，起到的应该是一种休息的作用。某些景点由于人流量大、服务不周等缺点，大大增加了游客数量在长时间内下降的可能性。我们除了必要的场景代入感，还应该兼具某些略显安静的观赏活动。例如在福建、江浙一带流行的功夫茶表演，这就给游客一种审美的享受。如果我们能够把一些酿造酒的技艺在现场重现给游客，让一些专业的品酒师举办知识普及的讲座，让游客有观赏和审美的体验，这无疑就带给了游客娱乐感和享受感.\newline 　　 最后，酒文化对于旅游事业还有经济发展的作用。除景点门票之外，旅游本身就会带动周边产业的发展，例如住宿、餐饮等服务性行业。另外，由于我们的这种弘扬酒文化的主张，会有大批的游客愿意到名酒原产地参观，购买纪念品。这又是一大收入。过去的种种事实证明，我们不能仅仅依靠门票来发展旅游业，这种单一的旅游经济来源必然会带来一些问题，从宏观上来看，这种方式不利于周边产业的发展，过于单一，时间一长，经济增长就会疲软。从长远来看，我们更需要形成产业链，扩大和挖掘这一产业的经济潜力。形成多种服务行业、甚至第一、第二产业的结合，全方位，多层次，宽领域地促进经济发展。}}
\section{感想和总结}
\paragraph{\textnormal{\newline 　　笔者来自北方，之前对贵州并没有深刻的了解。这种了解可以简单地形容为一种地域特点，比如内蒙古的草原，山西的黄土高坡等。而在此求学的将近一年的时间里，虽然去的地方并不算多，我还是逐渐了解到了贵州这一赋有民族特色的文化大省。所以，笔者更习惯在此文中用“我省”来代替贵州，因为这里让我产生了一种文化的亲近感。我们常说：“民族的，就是世界的”。这种民族文化无疑是我国甚至世界各民族文化之中闪耀的一支。\newline 　　自1915年万国博览会至今，已有一百又一年，自茅台酒走向世界也有一百多个年头了。茅台真正吸引大众的，恐怕不只是它的酱香风味。对于不是那么了解风味的普通民众来说，打动他们的，是那份几千年坚持古法工艺制造的那份匠人之心。茅台1400多种风味成分中，少了哪一道工序，都会使味道大打折扣，是很难在后期通过勾调出来的。笔者虽没能有机会品尝其中的芳香，但那份工人们对酿酒的传承和文化的沉淀却已在心中留下了独特的，难以抹去的馥郁之香。这就是匠心独具的魅力。\newline 　　在中国几千年的灿烂文明中，酒文化确是其中独特的一种文化。不论是民俗活动，还是酒文化对政治经济的影响作用，无疑是劳动人民智慧的结晶。“何以解忧，唯有杜康。”酒文化担得起这份几千年文化的担子，更经得起未来几千年的考验。酒，确能载得动许多愁！\newline 　　（学生才疏学浅，对酒文化及贵州旅游发展虽有自己的看法，但总体上来说还是知之甚少。如有不当之处，还望先生雅正。）
}}

\clearpage
%参考文献
\begin{thebibliography}{5}
    \zihao{-4}
    \bibitem{history} \zihao{-4}\textnormal{吴慧群,王仕佐. 利用“茅台酒”品牌优势 打造中国第一酒镇——关于茅台镇旅游业发展的思考[J]. 酿酒科技,2007,09:115-119.}
    \bibitem{culture} \zihao{-4}\textnormal{陈少波. 酒文化与贵州旅游[J]. 贵阳师专学报(社会科学版),1995,03:81-84.}
    \bibitem{function}\zihao{-4}\textnormal{杨之昉.从酒器谈唐代饮酒习俗[J].文博,2005,03:38-45.}
    \bibitem{react} \zihao{-4}\textnormal{吴正光. 贵州酒文化[J]. 当代贵州,2000,03:40-42.}
    \bibitem{history} \zihao{-4}\textnormal{王仕佐,邓咏梅,黄平. 略论贵州酒文化及旅游功能[J]. 酿酒科技,2003,05:84-87.}
\end{thebibliography}

\end{document}
